\paragraph{Area 3 – Inserimento delle prenotazioni}~\\
\subsubsection{[IU3.1] Modulo di prenotazione}
Nel modulo di prenotazione aula e docente si scelgono da elenchi di quelli registrati (non si digitano), il tipo di attività da un elenco, la data da un calendario. Un'opzione ``Ricorrenza settimanale'' abilita il campo della data di termine.\\
Priorità: SHOULD\\Fonte: Proposta del gruppo

\subsubsection{[IU3.2] Ricerca aule disponibili}
La ricerca delle aule disponibili ha i campi data, ora di inizio e ora di fine e mostra le aule libere in una tabella (codice, nome, edificio, piano, capienza), con il comando Prenota per l'aula selezionata (FI3.7).\\
Priorità: SHOULD\\Fonte: Traccia

\subsubsection{[IU3.3] Finestra dei conflitti}
In caso di conflitti in una serie il sistema mostra una finestra con: il numero di occorrenze in conflitto sul totale; per ogni occorrenza in conflitto, la data, il motivo (aula occupata o docente impegnato) e la prenotazione esistente con cui confligge; i comandi ``Salva solo le date libere'', ``Modifica dati'' e ``Annulla''. ``Salva solo le date libere'' è disabilitato se tutte le occorrenze sono in conflitto.\\
Priorità: MUST\\Fonte: C5
